\documentclass[a4paper]{article}

\usepackage{graphicx}
\usepackage{amsmath,amssymb}
\usepackage{minted}
\usepackage{multirow}
\usepackage{float}
\usepackage{algorithm}
\usepackage{algpseudocode}
\usepackage{todonotes}
\usepackage{forest}
\usepackage{xstring}
\usepackage{xcolor}
\usepackage[most]{tcolorbox}
\usepackage{xparse}
\usepackage{natbib}

\title{Repository reference navigation and Navigator selection fix}
\date{2026-09-04}
\author{Mohammad Rahmani}

\begin{document}
\maketitle

\section{Problem}
Local LaTeX references written with the portable \texttt{/repo/...} prefix were
mistaken for literal operating-system absolute paths. As a result, Ctrl+B could
not resolve those references to the Dropbox repository and therefore could not
open referenced LaTeX/YAML/text files in Mind Mirror or delegate PDFs/images to
the system application. In addition, switching the active editor tab could leave
the corresponding Navigator row without a visible selection highlight.

\section{Fix}
\begin{itemize}
   \item \texttt{WorkspaceReferenceResolver} now recognizes \texttt{/repo} as a repository alias and maps it through the configured workspace repository root before normal absolute-path handling.
   \item Existing command-specific extension inference remains active after alias resolution, including \texttt{includegraphics} targets without an explicit suffix.
   \item Ctrl+B continues to use the shared reference-opening path, so editable text files open inside Mind Mirror while binary/image/PDF files continue through the operating-system application association.
   \item Added \texttt{NavigatorActiveSelectionController} to make the active editor document the authoritative visible Navigator selection whenever the active tab changes.
\end{itemize}

\section{Regression coverage}
Tests cover direct \texttt{/repo} resolution, extension inference for graphics,
\texttt{href} occurrence resolution used by Ctrl+B, and installation of the
active-Navigator selection controller in the real enhanced window.

\section{Repository hygiene}
The implementation uses canonical source files and a dedicated controller class.
No patch, hotfix, backup implementation, or persistent temporary file is part of
the fix.

\end{document}
